The problem addressed in this project is the automated, non-invasive prediction of multiple tumor biomarkers in patients diagnosed with non-small cell lung cancer (NSCLC) using medical imaging data. In current clinical practice, biomarker evaluation, such as immunohistochemical (IHC) and genomic markers, is typically performed through tissue biopsies, which are invasive, expensive, and sometimes unfeasible for patients with advanced or inoperable disease. This project aims to develop an artificial intelligence (AI) model that can predict several clinically relevant tumor biomarkers directly from imaging data, using a hybrid approach that combines radiomics, deep learning, and patient-level demographic information.

This problem is of major clinical importance because accurate identification of tumor biomarkers guides personalized treatment strategies, including targeted therapies and immunotherapy. A reliable, non-invasive AI system would reduce the dependence on tissue biopsies, minimize patient risk, and enable faster, more accessible screening for actionable molecular markers. Such a tool could play a transformative role in precision oncology by providing physicians with immediate diagnostic insights before biopsy results are available.

The intended users of this AI system are radiologists, oncologists, and clinical researchers, who can use the tool as a decision-support system to complement manual image interpretation and laboratory tests. It is designed to assist in both diagnostic and treatment planning workflows by highlighting potential biomarker profiles based on imaging phenotypes.

The dataset used in this study consists of medical imaging data from 101 patients diagnosed with NSCLC, each having paired CT (Computed Tomography) and PET (Positron Emission Tomography) scans, with an average slice thickness of approximately 3.8 mm for CT and 3.3 mm for PET. All patients also underwent tissue biopsy, which provides the ground-truth labels for the corresponding tumor biomarkers used in training and validation. The cohort includes 42 female and 59 male patients, and the imaging data are stored in standard DICOM format. The paired multimodal volumes, optionally supplemented with basic clinical metadata such as patient age or sex, are processed through an AI pipeline that performs image preprocessing, feature extraction (radiomics and deep features), and predictive modeling. The output of the system is a set of predictions indicating the likelihood of the presence or expression level of multiple tumor biomarkers, such as TTF1, NAPSIN, and CK7, each represented as a confidence score or classification label. \cite{Shao2022-yo} This comprehensive dataset, together with biopsy-confirmed biomarker status, provides a solid foundation for developing and evaluating fusion-based deep learning models, allowing the system to learn complementary metabolic and anatomical signatures relevant for biomarker prediction.


To explore the impact of multimodal fusion strategies, we will develop three separate AI models that all rely on DenseNet-based feature extractors, but differ in how PET and CT information is combined. The first model will use early fusion, merging the modalities at the input or initial feature-map level. The second will use intermediate (late) fusion, combining deep features only after several DenseNet blocks. The third model will also follow a late fusion design, but with shared convolutional weights between the PET and CT pathways, ensuring that both modalities learn synchronized feature representations. This setup allows a systematic comparison of fusion mechanisms while maintaining architectural consistency across models.~\cite{Da-Ano2024-oo} In all three cases, the fused deep features will be further combined with radiomic features at the decision stage, enabling a hybrid prediction mechanism that leverages both learned representations and quantitative imaging biomarkers.~\cite{Li2024-np}


Model performance is evaluated quantitatively using standard classification and regression metrics, depending on the biomarker type. The main metrics include the AUC, sensitivity, specificity, and F1-score. Cross-validation on the training data and testing on independent validation cohorts are applied to ensure the robustness, generalizability, and clinical reliability of the model. 

Overall, the project seeks to create a clinically interpretable, reproducible, and data-driven virtual biopsy system capable of predicting multiple tumor biomarkers from PET/CT imaging, thus contributing to the ongoing shift toward precision medicine in lung cancer management.
